
\begin{frame}{Opérateur =  itérateur}

Tout opérateur est implanté sous forme d'un \alert{itérateur}.
Trois fonctions:

\begin{redbullet}
  \item \texttt{open}\,: initialise les ressources et 
            positionne le curseur\,;
   
  \item \texttt{next}\,: ramène l'enregistrement courant\;
          se place sur l'enregistrement suivant\,;

  \item \texttt{close}\,: libère les ressources\,;
\end{redbullet}

\vfill

\'Echanges: 

\begin{blocgauche}
\begin{redbullet}
  \item Un itérateur \alert{consomme} des nuplets d'autres itérateurs \alert{source}.
  \item Un itérateur \alert{produit} des nuplets pour un autre
     itérateur (ou pour l'application).
\end{redbullet}  
\end{blocgauche}
    
    
\end{frame}


\begin{frame}[fragile]{Exemple: parcours d'index (\texttt{IndexScan})}

Rappel: index = arbre B. 

\vfill

\begin{redbullet} 
  \item  Pendant le \texttt{open()}: parcours de la racine vers la feuille.
  \item À chaque appel à \texttt{next()}: parcours en séquence des feuilles.  
\end{redbullet}


\figSlideWithSize{iter-indexScan}{10cm}

\smallskip
\begin{blocgauche}
\alert{Efficacité}: très efficace, quelques lectures logiques (index en mémoire)
\end{blocgauche}
\end{frame}

%%%%
\begin{frame}[fragile]{Accès par adresse: \texttt{DirectAccess}}

\begin{redbullet}
  \item Pendant le \texttt{open()}: rien à faire. 
  \item À chaque appel à \texttt{next()}: on reçoit une adresse, on produit un nuplet. 
\end{redbullet}  

\vfill

\figSlideWithSize{planEx-DirectAccess}{10cm}

\vfill

Très efficace: un accès bloc, souvent en mémoire.
\end{frame}

\begin{frame}{Plan d'exécution}

Un plan d'exécution connecte les opérateurs. Ici, recherche avec index.

\vfill

\figSlideWithSize{planEx-pipeline}{12cm}

\vfill

Le pipelinage reste complet.

\end{frame}



\begin{frame}[fragile]{Un exemple de base}

Nous allons étudier les plans permettant d'exécuter les
requêtes \alert{mono-table}.

\begin{minted}[frame=leftline,
               baselinestretch=.9,
               fontsize=\footnotesize,
               framesep=2mm]{sql}
select a1, a2, ..., an 
from T 
where condition
\end{minted}

\vfill

De quels opérateurs a-t-on besoin\,?

\begin{redbullet}
 \item ~ [\texttt{FullScan}]: parcours séquentiel de la table (déjà vu).
 \item ~ [\texttt{IndexScan}]: parcours d'un index (si disponible).
 \item ~ [\texttt{DirectAccess}]: accès \alert{par adresse} à un nuplet.
 \item ~ [\texttt{Filter}]: test de la condition.
\end{redbullet}

\vfill

Nous obtenons \alert{deux} plans d'exécution possibles.
\end{frame}

\begin{frame}[fragile]{Premier plan d'exécution: sans index}

\begin{minted}[frame=leftline,
               baselinestretch=.9,
               fontsize=\footnotesize,
               framesep=2mm]{sql}
select titre from Film where genre='SF' and annee = 2014
\end{minted}

\vfill


\figSlideWithSize{planEx-mono-scan}{11cm}

\vfill
\begin{blocgauche}
N'utilise pas d'index
\end{blocgauche}
\end{frame}

\begin{frame}[fragile]{Second plan d'exécution: avec index}

\begin{minted}[frame=leftline,
               baselinestretch=.9,
               fontsize=\footnotesize,
               framesep=2mm]{sql}
select titre from Film where genre='SF' and annee = 2014
\end{minted}

\vfill

\figSlideWithSize{planEx-mono-index}{12cm}

\vfill
\begin{blocgauche}
Utilise un index sur l'année.
\end{blocgauche}
\end{frame}


\begin{frame}[fragile]{Index ou pas index?}

Recherche des nuplets entre 9 et 13, avec 3 blocs en mémoire.

\medskip


\only<1-1|handout:1-1>{
\figSlideWithSize{interBtree}{10cm}
}
\only<2-2|handout:2-2>{
\figSlideWithSize{interBtree-1}{10cm}
}
\only<3-3|handout:3-3>{
\figSlideWithSize{interBtree-2}{10cm}
}
\only<4-4|handout:4-4>{
\figSlideWithSize{interBtree-3}{10cm}
}
\only<5-5|handout:5-5>{
\figSlideWithSize{interBtree-4}{10cm}
}
\only<6-6|handout:6-6>{
\figSlideWithSize{interBtree-5}{10cm}
}
\only<7-7|handout:7-7>{
\vfill
Avec l'index, il faut lire 5 blocs! Parcours séquentiel bien préférable.
\vfill
\begin{blocgauche}
Un cas extrême!  \alert{En pratique}: le SGBD décide en fonction des statistiques
et des ressources disponibles.
\end{blocgauche}
}

\end{frame}



%%%%
\begin{frame}[fragile]{Opérateur de sélection}

Très simple: filtre les nuplets fournis par la source.

\begin{minted}[frame=leftline,
               baselinestretch=.9,
               fontsize=\footnotesize,
               framesep=2mm]{bash}
    function nextFilter
    {
       # On prend un nuplet de la source
      $nuplet = $source.next();
      # On continue tant que la condition n'est pas satisfaite, 
      # ou la source parcourue
      while ($nuplet != null  and $nuplet.test($C) = false)
      do
        $nuplet = $source.next();        
      done
      
      return $nuplet;
    }
\end{minted}

\end{frame}



\begin{frame}[fragile]{Pour compléter}

Un plan à exécuter avec index.

\begin{minted}[frame=leftline,
               baselinestretch=.9,
               fontsize=\footnotesize,
               framesep=2mm]{sql}
    select * from   Film
    where  idFilm = 20 
    and    titre = 'Vertigo'
\end{minted}

\vfill

Un plan à exécuter sans index.

\begin{minted}[frame=leftline,
               baselinestretch=.9,
               fontsize=\footnotesize,
               framesep=2mm]{sql}
    select * from   Film
    where  idFilm = 20 
    or    titre = 'Vertigo'
\end{minted}

\end{frame}


\begin{frame}{Résumé: plans pour requêtes mono-table}

Permier aperçu de l'optimisation

\vfill

\begin{redbullet}
  \item Le système a le choix entre plusieurs plans possibles.
  \item Distinguer le plus efficace n'est pas toujours trivial.
  \item Le choix peut changer selon le contexte.
\end{redbullet}


\vfill

\only<2|handout:2>{
\vfill
\centerline{\alert{Merci!}}
}


\end{frame}
